\section{引言}

给定实体对及其图上下文，关系预测需要识别相应的关系标签。GRIP 将图结构融入语言模型参数，通过生成式问答完成这一任务 \citep{feng2025grip}。候选级对比学习让目标关系与替代关系直接竞争，而候选的选取决定模型接收到什么关系辨别信号。因此，候选构造是连接图结构与关系排序目标的关键环节。

困难候选的构造需要同时处理结构相关性、模型混淆与事实不确定性。训练图的关系统计描述不同关系的结构接近程度，模型分数则揭示已有的关系混淆。知识图谱的不完整性进一步要求区分候选的学习用途与事实真值：未标注关系在开放世界中保留未判定状态。本文研究如何将结构与模型混淆信号共同用于关系候选构造，并将训练已知关系过滤和来源追溯纳入同一流程。

为此，我们提出结构反事实对比学习（\scc）。\scc{} 以关系结构画像刻画候选在训练图中的统计接近程度，以冻结的 generation-only 基线 \bone{} 的分数违例度刻画模型混淆，两种信号共同决定带均匀下限的采样分布。训练已知关系过滤与固定候选库明确候选来源，使每个候选能够关联到其构造记录。\scc{} 接入 GRIP Stage~2，联合优化候选级 InfoNCE 与原始生成目标，沿用 Stage~1 和原始提示。本文将结构反事实定义为固定实体对上的关系标签替换，图~\ref{fig:mechanism} 展示其构造与学习流程。

在 NELL23K 的 4,943 条测试问题上，\scc{} 的候选重排 H@1 为 91.52\%，MRR 为 95.27\%。在相同候选列表的前 200 条配对问题上，其 MRR 比 \bone{} 高 1.60 个百分点，H@1 高 2.50 个百分点。实验节分别展示完整测试集的排序水平和配对子集的逐题预测变化，并报告直接解码结果、候选来源检查及分阶段成本。评估设置与比较范围见第~\ref{sec:experiments} 节。

本文面向固定实体对上的关系候选排序，贡献包括三个方面：定义结构画像与冻结模型违例度联合驱动的训练候选分布；构建包含训练已知关系过滤、固定候选库和逐候选来源记录的学习流程；在 NELL23K 展示全量候选重排结果与同候选前 200 题的配对分析。
